%% ============================================================
%%  Aula 13 — Busca em árvore de Monte Carlo (MCTS) e AlphaZero
%%  Versão sucinta: roteiro do apresentador (poucas palavras,
%%  mesmas definições e figuras do aula13.tex)
%%  Adaptação em pt-BR do CS234 (Stanford, Emma Brunskill), Lecture 13
%%  Compilar: latexmk -xelatex aula13.tex
%% ============================================================
\documentclass[aspectratio=169, 11pt]{beamer}
%% .sty do projeto na raiz do repositório (../)
\makeatletter
\def\input@path{{./}{../}}
\makeatother


\usetheme{AprendizadoRL}      % ANTES do babel — fontspec precisa vir primeiro
\usepackage{babel}
\babelprovide[import, main]{brazilian}
\usepackage{booktabs}
\usepackage{pgfplots}
\pgfplotsset{compat=1.18}

\title[Aula 13 — MCTS e AlphaZero]{Aula 13: Busca em árvore de Monte Carlo\\ e a família AlphaZero}
\subtitle[Aprendizagem por Reforço]{Busca baseada em simulação \textbullet\ UCT \textbullet\ AlphaGo, AlphaZero e MuZero}
\author[Prof. Marcelo K. Albertini]{Prof. Marcelo Keese Albertini \\
  \footnotesize Universidade Federal de Uberlândia \\
  \footnotesize adaptado de Emma Brunskill, CS234 (Stanford)}
\institute{Universidade Federal de Uberlândia}
\date{2026}

% ---- Nota discreta: perguntas ficam para o professor conduzir em aula
\newcommand{\paradiscussao}{\smallskip
  {\footnotesize\color{rlCinza}\itshape Pergunta para discussão conduzida em aula.}}

% ---- Macros locais desta aula
\newcommand{\Mmod}{\mathcal{M}}                 % modelo / simulador
\newcommand{\Nv}[1]{N(#1)}                      % contagem de visitas
\newcommand{\Qh}{\widehat{Q}}                   % estimativa empírica
\newcommand{\Vh}{\widehat{V}}
\newcommand{\ret}{G}                            % retorno
\newcommand{\raiz}{s_{\mathrm{raiz}}}
\newcommand{\folha}{s_{\mathrm{folha}}}

% ---- Estilo de caixa TikZ (o tema Beamer não define `caixa`, só o handout)
\tikzset{caixa/.style={draw, thick, rounded corners=2pt, align=center,
                       font=\scriptsize}}

% ---- Árvore genérica reutilizável para as fases do MCTS
\newcommand{\mctsarvore}[1]{%
\begin{tikzpicture}[baseline=(r.north), scale=0.88, transform shape]
  \tikzset{no/.style={circle, draw=rlAzul, fill=rlAzul!10, thick,
                      minimum size=4.6mm, inner sep=0pt}}
  \node[no] (r) at (0,0) {};
  \node[no] (a) at (-0.95,-0.85) {};
  \node[no] (b) at (0.95,-0.85) {};
  \node[no] (c) at (-1.45,-1.7) {};
  \node[no] (d) at (-0.45,-1.7) {};
  \draw[trans] (r)--(a);  \draw[trans] (r)--(b);
  \draw[trans] (a)--(c);  \draw[trans] (a)--(d);
  #1
\end{tikzpicture}}

\begin{document}

\begin{frame}[plain, noframenumbering]\titlepage\end{frame}

%% ------------------------------------------------------------
\begin{frame}{Plano de hoje}
  \begin{enumerate}
    \item \textbf{Busca baseada em simulação} — decidir bem \emph{agora}, em vez
          de calcular política para \emph{todo} o espaço de estados.
    \item \textbf{Busca de Monte Carlo simples e árvores expectimax} — o que se
          ganha; enumeração explode.
    \item \textbf{MCTS e UCT} — quatro fases, o bandit em cada nó,
          garantias e limitações.
    \item \textbf{AlphaGo $\to$ AlphaGo Zero $\to$ AlphaZero $\to$ MuZero} — \emph{a busca é um operador de melhoria de política}.
  \end{enumerate}

  \vfill
  \begin{ideiabox}
    Pergunta de fundo: com um simulador do mundo, quanto vale \emph{pensar mais} antes de agir --- e como converter isso em aprendizado permanente?
  \end{ideiabox}
\end{frame}

%% ------------------------------------------------------------
\begin{frame}{Onde estamos no curso}{Da política global ao planejamento local}
  \centering
  \begin{tikzpicture}[node distance=3mm]
    \node[caixa, draw=rlAzul!45, fill=rlAzul!6, text width=2.7cm,
          rounded corners=2pt, align=center, font=\scriptsize, inner sep=3pt]
          (a) {\textbf{Aulas 2--4}\\ Programação dinâmica, MC, TD, Q-learning\\[2pt]
               \color{rlCinza}política para \emph{todo} $\gS$};
    \node[caixa, draw=rlAzul!45, fill=rlAzul!6, text width=2.7cm, right=6mm of a,
          rounded corners=2pt, align=center, font=\scriptsize, inner sep=3pt]
          (b) {\textbf{Aulas 5--8}\\ Gradientes de política, imitação, RLHF\\[2pt]
               \color{rlCinza}política parametrizada};
    \node[caixa, draw=rlAzul!45, fill=rlAzul!6, text width=2.7cm, right=6mm of b,
          rounded corners=2pt, align=center, font=\scriptsize, inner sep=3pt]
          (c) {\textbf{Aulas 9--12}\\ Exploração, bandits, RL eficiente em amostras\\[2pt]
               \color{rlCinza}quanto custa aprender};
    \node[caixa, draw=rlLaranja, fill=rlLaranja!10, text width=2.7cm, right=6mm of c,
          rounded corners=2pt, align=center, font=\scriptsize, inner sep=3pt]
          (d) {\textbf{Aula 13}\\ Busca por simulação, MCTS, AlphaZero\\[2pt]
               \color{rlLaranja}\bfseries decidir bem \emph{agora}};
    \draw[trans] (a)--(b); \draw[trans] (b)--(c); \draw[trans destac] (c)--(d);
  \end{tikzpicture}

  \vfill
  \begin{itemize}
    \item Recurso novo: \textbf{um modelo do qual se pode amostrar}.
    \item Além de jogos: \textbf{gastar computação em tempo de decisão} (\emph{test-time compute}).
  \end{itemize}
\end{frame}

\section{Busca baseada em simulação}

%% ------------------------------------------------------------
\begin{frame}{Calcular a ação apenas para o estado atual}
  Até aqui: objetos \emph{globais} — $\Vstar(s)$ para todo $s$, ou $\pol_\theta(a\mid s)$ para toda entrada.

  \pause
  \begin{ideiabox}
    \textbf{Ideia central da aula.} Basta uma boa ação \emph{no estado em que estou}. Se sobra computação, gaste-a localmente, no subproblema que começa agora.
  \end{ideiabox}

  \pause
  \begin{itemize}
    \item Agente em $s_t$ (estado \emph{real}); simulações internas a partir de $s_t$ com um modelo $\Mmod$.
    \item Só a raiz importa: busca produz $a_t$; executa no mundo real, observa $s_{t+1}$, recomeça.
    \item \textbf{Planejamento em tempo de decisão} --- oposto do \emph{em segundo plano} (iteração de valor).
  \end{itemize}
\end{frame}

%% ------------------------------------------------------------
\begin{frame}{O que exatamente é ``ter um modelo''}
  \begin{defbox}{Modelo generativo (simulador)}
    Caixa-preta $\Mmod$ que, dado $(s,a)$, devolve uma \emph{amostra}
    $s' \sim \tran{\cdot}{s,a}$ e $r \sim \rec(\cdot\mid s,a)$. Não é preciso
    conhecer as probabilidades, nem varrer todos os sucessores.
  \end{defbox}

  \pause
  \begin{itemize}
    \item Jogos de tabuleiro: simulador \emph{perfeito e gratuito} — as regras são o modelo.
    \item Robótica e saúde: aprendido e imperfeito --- o erro acumula ao longo do horizonte simulado.
    \item MCTS é agnóstico: exige apenas amostras.
  \end{itemize}

  \pause
  \begin{alertabox}
    Bom para \emph{um passo} pode ser inútil para \emph{cinquenta}. A busca vale até onde o modelo é fiel.
  \end{alertabox}
\end{frame}

%% ------------------------------------------------------------
\begin{frame}{Busca de Monte Carlo simples}
  Dados um modelo $\Mmod$ e uma \emph{política de simulação} $\pol$:

  \vspace{-10pt}
  \begin{algbox}{Busca de Monte Carlo simples}
    \begin{itemize}\setlength{\itemsep}{0.1ex}
      \item Para cada ação $a \in \gA$, simule $K$ episódios a partir do estado
            real $s_t$:
            $\{s_t, a, r^k_{t+1}, s^k_{t+1}, \dots, s^k_T\}_{k=1}^{K}
             \sim \Mmod, \pol$, e estime
            \[
              \Qh(s_t,a) \;=\; \frac{1}{K}\sum_{k=1}^{K} \ret^k_t
              \;\xrightarrow[K\to\infty]{}\; \Qpi(s_t,a).
            \]
      \item Execute no mundo real $\displaystyle a_t = \argmax_{a\in\gA} \Qh(s_t,a)$.
    \end{itemize}
  \end{algbox}

  \pause
  \begin{ideiabox}
    Isso é \textbf{um passo de melhoria de política} sobre $\pol$, feito por amostragem, apenas no estado $s_t$.
  \end{ideiabox}
\end{frame}

%% ------------------------------------------------------------
\begin{frame}{Por que é exatamente um passo de melhoria}
  Seja $\pol'$ a política que age gulosamente sobre $\Qpi$ em $s_t$ e segue
  $\pol$ depois. Pelo teorema de melhoria de política (Aula 3),
  \[
    V^{\pol'}(s_t) \;=\; \max_{a} \Qpi(s_t,a) \;\ge\; \Vpi(s_t).
  \]

  \pause
  \begin{itemize}
    \item A busca \emph{estima} o $\max_a \Qpi(s_t,a)$ por amostragem: com $K\to\infty$ recuperamos $\pol'$ exatamente.
    \item Teto: nunca melhor que $\pol'$. Se $\pol$ é ruim, um passo continua sendo pouco.
    \item Episódios abaixo do primeiro nível: \emph{desperdiçados} como decisão — só avaliam $\pol$.
  \end{itemize}

  \pause
  \begin{alertabox}
    $\Qh$ é média de retornos --- variância alta. $\argmax$ de estimativas ruidosas: \textbf{viés de maximização} — escolhe-se a ação \emph{com mais sorte}.
  \end{alertabox}
\end{frame}

%% ------------------------------------------------------------
\begin{frame}{Podemos fazer melhor que um passo?}
  \begin{itemize}
    \item Sim: \emph{repetidamente} o backup de Bellman ao longo da simulação --- buscar em profundidade, maximizando em cada nível.
    \item Objeto ideal: \textbf{árvore expectimax} — alterna $\max$ sobre ações e $\E$ sobre transições.
  \end{itemize}

  \pause
  \begin{teobox}{Backup expectimax até o horizonte $H$}
    \[
      Q_h(s,a) \;=\; \rec(s,a) \;+\; \gamma \sum_{s'} \tran{s'}{s,a}\,
      \underbrace{\max_{a'} Q_{h-1}(s',a')}_{V_{h-1}(s')},
      \qquad Q_0 \equiv 0 .
    \]
    Na raiz, $\argmax_a Q_H(s_t,a)$ é ótima para o sub-MDP de $H$ passos que
    começa em $s_t$.
  \end{teobox}

  \pause
  Não é preciso resolver o MDP inteiro --- basta o \emph{sub-MDP} que começa agora.
\end{frame}

%% ------------------------------------------------------------
\begin{frame}{A árvore expectimax}
  \centering
  \begin{tikzpicture}[node distance=6mm]
    \node[estado destacado] (s) at (0,0) {$s_t$};
    \node[acaon] (a1) at (-2.4,-1.1) {};
    \node[acaon] (a2) at (0,-1.1) {};
    \node[acaon] (a3) at (2.4,-1.1) {};
    \foreach \n/\lb in {a1/$a_1$, a2/$a_2$, a3/$a_3$}
      \draw[trans destac] (s) -- node[left, font=\scriptsize, color=rlLaranja, inner sep=1pt] {\lb} (\n);
    \foreach \p/\x in {a1/-2.4, a2/0, a3/2.4}{
      \node[estado, minimum size=6mm, font=\tiny] (l\p) at (\x-0.75,-2.2) {};
      \node[estado, minimum size=6mm, font=\tiny] (r\p) at (\x+0.75,-2.2) {};
      \draw[trans] (\p) -- node[probl] {\tiny $p$} (l\p);
      \draw[trans] (\p) -- node[probl] {\tiny $1{-}p$} (r\p);}
    \node[right=2mm of s, font=\scriptsize, color=rlLaranja] {\bfseries $\max$};
    \node[font=\scriptsize, color=rlCinza] at (3.9,-1.1) {$\E$ sobre $\tran{s'}{s,a}$};
    \node[font=\scriptsize, color=rlCinza] at (0,-2.9) {$\vdots$ até profundidade $H$};
  \end{tikzpicture}

  \vfill
  \begin{itemize}
    \item Nós circulares: estados (\emph{esperança} sobre a dinâmica).
    \item Nós sólidos: ações (\emph{maximização}).
    \item Soma zero: nível do oponente vira $\min$ — \emph{minimax}.
  \end{itemize}
\end{frame}

%% ------------------------------------------------------------
\begin{frame}{O problema: a árvore explode}
  \begin{teobox}{Custo da enumeração}
    Uma árvore expectimax completa de profundidade $H$ tem
    $\;\bigl(\abs{\gS}\,\abs{\gA}\bigr)^{H}\;$ nós.
  \end{teobox}

  \pause
  \begin{itemize}
    \item Xadrez: fator de ramificação $b \approx 35$, partida típica $d \approx 80$.
    \item \emph{Go} $19\times19$: $b \approx 250$, $d \approx 150$. $250^{150}$ é maior que qualquer número físico relevante.
    \item Poda alfa-beta (minimax determinístico): expoente cai pela metade --- ainda impossível.
  \end{itemize}

  \pause
  \begin{ideiabox}
    Duas saídas --- o MCTS usa \textbf{as duas}: \textbf{amostrar} em vez de enumerar (largura) e \textbf{truncar} com estimativa de valor (profundidade).
  \end{ideiabox}
\end{frame}

%% ------------------------------------------------------------
\begin{frame}{Para discutir}
  \begin{quizbox}{}
    Com orçamento fixo de $N$ chamadas ao simulador, vale mais distribuir as
    simulações igualmente entre as $\abs{\gA}$ ações da raiz, ou concentrá-las
    nas que \emph{até agora} parecem boas? Em que regime a resposta muda?
  \end{quizbox}
  \paradiscussao

  \pause
  \vfill
  \begin{quizbox}{}
    Por que a poda alfa-beta, tão eficaz no xadrez dos anos 1990, não resolveu
    o \emph{Go}? Que propriedade do problema ela exige e o \emph{Go} não oferece?
  \end{quizbox}
  \paradiscussao
\end{frame}

\section{MCTS e a regra UCT}

%% ------------------------------------------------------------
\begin{frame}{Busca em árvore de Monte Carlo (MCTS)}
  \begin{algbox}{Esqueleto do MCTS}
    Dado um modelo $\Mmod$ e um orçamento de $K$ simulações:
    \begin{enumerate}\setlength{\itemsep}{0.25ex}
      \item Construa uma árvore de busca enraizada no estado atual $s_t$.
      \item Repita $K$ vezes: percorra a árvore \emph{amostrando} ações e
            estados sucessores, expanda-a em um ponto, estime o valor da nova
            folha e propague a estimativa de volta até a raiz.
      \item Ao fim do orçamento, escolha a ação real
            $a_t = \argmax_{a\in\gA} f\bigl(\raiz, a\bigr)$, onde $f$ é o valor
            estimado ou a contagem de visitas.
    \end{enumerate}
  \end{algbox}

  \pause
  \begin{itemize}
    \item Árvore \textbf{assimétrica}: fundo onde promissor, raso onde não — o oposto da busca em largura.
    \item \textbf{\emph{Anytime}}: interrompida a qualquer momento, devolve a melhor decisão até ali.
  \end{itemize}
\end{frame}

%% ------------------------------------------------------------
\begin{frame}{As quatro fases de uma simulação}
  \centering
  \begin{tabular}{@{}cccc@{}}
    \mctsarvore{\draw[trans destac] (r)--(a); \draw[trans destac] (a)--(d);
                \node[font=\tiny, color=rlLaranja, below=1mm of d] {?};}
    &
    \mctsarvore{\draw[trans destac] (r)--(a); \draw[trans destac] (a)--(d);
                \node[circle, draw=rlLaranja, fill=rlLaranja!25, thick,
                      minimum size=4.6mm, inner sep=0pt] (e) at (-0.45,-2.55) {};
                \draw[trans destac] (d)--(e);}
    &
    \mctsarvore{\node[circle, draw=rlLaranja, fill=rlLaranja!25, thick,
                      minimum size=4.6mm, inner sep=0pt] (e) at (-0.45,-2.55) {};
                \draw[trans] (d)--(e);
                \draw[rlVerde!70!black, thick, decorate,
                      decoration={snake, amplitude=0.5mm, segment length=2.2mm},
                      -{Stealth[length=2mm]}] (e) -- ++(0,-0.85);
                \node[font=\tiny, color=rlVerde!60!black] at (0.6,-3.4) {$\ret$};}
    &
    \mctsarvore{\node[circle, draw=rlLaranja, fill=rlLaranja!25, thick,
                      minimum size=4.6mm, inner sep=0pt] (e) at (-0.45,-2.55) {};
                \draw[trans] (d)--(e);
                \draw[{Stealth[length=2mm]}-, very thick, rlVerde!70!black] (d)--(e);
                \draw[{Stealth[length=2mm]}-, very thick, rlVerde!70!black] (a)--(d);
                \draw[{Stealth[length=2mm]}-, very thick, rlVerde!70!black] (r)--(a);}
    \\[0.4em]
    {\scriptsize\bfseries\color{rlAzul} 1. Seleção} &
    {\scriptsize\bfseries\color{rlAzul} 2. Expansão} &
    {\scriptsize\bfseries\color{rlAzul} 3. Simulação} &
    {\scriptsize\bfseries\color{rlAzul} 4. Retropropagação} \\[0.15em]
    {\tiny\color{rlCinza}\parbox{2.9cm}{\centering desce pela árvore com a
      \emph{política de árvore} (UCT)}} &
    {\tiny\color{rlCinza}\parbox{2.9cm}{\centering acrescenta um nó novo ao
      encontrar ação não tentada}} &
    {\tiny\color{rlCinza}\parbox{2.9cm}{\centering estima o valor da folha:
      \emph{rollout} ou rede de valor}} &
    {\tiny\color{rlCinza}\parbox{2.9cm}{\centering atualiza $N$ e $Q$ em todas
      as arestas do caminho}}
  \end{tabular}

  \vspace{-3pt}
  \begin{ideiabox}
    Dentro da árvore a política \emph{muda a cada simulação}. Fora dela, o \emph{rollout} usa uma política fixa e barata.
  \end{ideiabox}
\end{frame}

%% ------------------------------------------------------------
\begin{frame}{O que fica guardado em cada aresta}
  Para cada par $(i,a)$ --- nó $i$ da árvore, ação $a$ --- o algoritmo mantém:

  \vspace{0.4em}
  \centering
  \begin{tabular}{@{}lll@{}}
    \toprule
    \textbf{Estatística} & \textbf{Significado} & \textbf{Atualização} \\
    \midrule
    $\Nv{i,a}$ & vezes que $a$ foi escolhida em $i$ & $+1$ a cada visita \\
    $W(i,a)$   & soma dos retornos observados       & $+\,\ret$ \\
    $\Qh(i,a)$ & valor médio $W(i,a)/\Nv{i,a}$      & recalculado \\
    $\Nv{i}$   & $\sum_a \Nv{i,a}$, visitas ao nó   & $+1$ \\
    \bottomrule
  \end{tabular}

  \vfill
  \raggedright
  \begin{itemize}
    \item \emph{Médias de amostras} --- nada de matriz de transição, nada de varredura sobre $\gS$.
    \item Memória cresce com nós \emph{visitados}, não com $\abs{\gS}$ — por isso MCTS roda em jogos com $10^{170}$ posições.
  \end{itemize}
\end{frame}

%% ------------------------------------------------------------
\begin{frame}{Cada nó é um problema de bandit}
  Escolher qual ramo simular é o dilema exploração/aproveitamento da Aula 10:

  \begin{itemize}
    \item cada ação $a$ no nó $i$ é um \emph{braço};
    \item o ``ganho'' é o retorno $\ret$ da simulação que passa por ele;
    \item gastar nos braços bons, sem descartar cedo demais um braço que teve azar.
  \end{itemize}

  \pause
  \begin{teobox}{UCB1 (Auer, Cesa-Bianchi e Fischer, 2002) --- revisão}
    \[
      a_t \;=\; \argmax_{a}\; \Qh(a) \;+\; c\sqrt{\frac{\ln t}{N(a)}} .
    \]
    O segundo termo é a largura de um intervalo de confiança: encolhe como
    $1/\sqrt{N(a)}$ e cresce com o tempo, de modo que nenhum braço é abandonado
    para sempre.
  \end{teobox}

  \pause
  \textbf{UCT} = UCB aplicado a árvores: use UCB1 \emph{em cada nó}.
\end{frame}

%% ------------------------------------------------------------
\begin{frame}{UCT: \emph{Upper Confidence Trees}}
  \begin{teobox}{Regra de seleção UCT (Kocsis e Szepesvári, 2006)}
    No nó $i$, selecione
    \[
      a \;=\; \argmax_{a\in\gA(i)}\;
      \mdestaque{rlAzulClaro}{\Qh(i,a)}
      \;+\;
      \mdestaque{rlLaranja}{c\sqrt{\frac{\ln \Nv{i}}{\Nv{i,a}}}}
      \qquad
      \Qh(i,a) = \frac{1}{\Nv{i,a}}\sum_{k=1}^{\Nv{i,a}} \ret^k(i,a)
    \]
    com $\ret^k(i,a)$ o $k$-ésimo retorno a partir de $i$ após $a$ e
    $\Nv{i}=\sum_a \Nv{i,a}$.
  \end{teobox}

  \pause
  \begin{itemize}\setlength{\itemsep}{0.3ex}
    \item Ação nunca tentada tem $\Nv{i,a}=0$: bônus infinito --- visitada primeiro (a \emph{expansão}).
    \item $c$ depende da escala dos retornos: normalizar $\ret$ é padrão.
    \item Política de árvore \textbf{muda a cada simulação} --- melhoria de política \emph{durante} a busca.
  \end{itemize}
\end{frame}

%% ------------------------------------------------------------
\begin{frame}{Como o bônus decai}
  \centering
  \begin{tikzpicture}
    \begin{axis}[width=9.4cm, height=4.6cm,
      xlabel={visitas ao ramo $\Nv{i,a}$}, ylabel={bônus},
      xmin=1, xmax=200, ymin=0, ymax=2.4,
      grid=major, grid style={rlAzul!12},
      label style={font=\scriptsize, color=rlCinza},
      tick label style={font=\tiny, color=rlCinza},
      legend style={font=\tiny, draw=rlAzul!25, fill=white},
      legend pos=north east, axis lines=left,
      every axis plot/.append style={thick}]
      \addplot[rlLaranja, domain=1:200, samples=120]
        {sqrt(ln(1000)/x)};  \addlegendentry{$\Nv{i}=1000$}
      \addplot[rlAzulClaro, domain=1:200, samples=120]
        {sqrt(ln(200)/x)};   \addlegendentry{$\Nv{i}=200$}
      \addplot[rlTeal, domain=1:200, samples=120]
        {sqrt(ln(50)/x)};    \addlegendentry{$\Nv{i}=50$}
    \end{axis}
  \end{tikzpicture}

  \vfill
  \raggedright
  \begin{itemize}
    \item Cai como $\Nv{i,a}^{-1/2}$ --- rápido no começo, lento depois.
    \item $\sqrt{\ln \Nv{i}}$: dobrar o orçamento mal move a curva. MCTS \emph{concentra} o esforço.
    \item Visitas na raiz fortemente desbalanceadas --- é o que as torna informativas.
  \end{itemize}
\end{frame}

%% ------------------------------------------------------------
\begin{frame}{MCTS/UCT em pseudocódigo}
  {\scriptsize
  \begin{algbox}{UCT}
    \begin{itemize}\setlength{\itemsep}{0.15ex}
      \item[] \textbf{entrada:} $s_t$, modelo $\Mmod$, orçamento $K$, constante $c$;
              árvore $\gets$ raiz $\raiz = s_t$
      \item[] \textbf{para} $k = 1, \dots, K$:
      \item[] \quad $s \gets \raiz$; \; caminho $\gets [\,]$
      \item[] \quad \textbf{enquanto} $s$ está na árvore e todas as ações de $s$ já foram tentadas:
      \item[] \quad\quad $a \gets \argmax_{a} \Qh(s,a) + c\sqrt{\ln \Nv{s}/\Nv{s,a}}$
              \hfill\textit{\color{rlCinza}\scriptsize seleção}
      \item[] \quad\quad $(s', r) \sim \Mmod(s,a)$; \; anexe $(s,a,r)$ ao caminho; \; $s \gets s'$
      \item[] \quad \textbf{se} $s$ não é terminal: acrescente $s$ à árvore
              \hfill\textit{\color{rlCinza}\scriptsize expansão}
      \item[] \quad $\ret \gets$ \textsc{Avalia}$(s)$
              \hfill\textit{\color{rlCinza}\scriptsize simulação}
      \item[] \quad \textbf{para} $(s,a,r)$ no caminho, de trás para frente:
      \item[] \quad\quad $\ret \gets r + \gamma \ret$; \;
              $\Nv{s,a} \mathrel{+}= 1$; \;
              $W(s,a) \mathrel{+}= \ret$; \;
              $\Qh(s,a) \gets W(s,a)/\Nv{s,a}$
              \hfill\textit{\color{rlCinza}\scriptsize retropropagação}
      \item[] \textbf{devolva} $\argmax_{a} \Nv{\raiz,a}$
      \item[] \rule{0pt}{1.1em}\textsc{Avalia}$(s)$: \emph{rollout} até o fim
              (MCTS clássico) ou $v_\theta(s)$ de uma rede de valor (AlphaZero).
    \end{itemize}
  \end{algbox}}
\end{frame}

%% ------------------------------------------------------------
\begin{frame}{Detalhes que decidem se funciona}
  \begin{itemize}
    \item \textbf{Ação final:} $\argmax_a \Nv{\raiz,a}$, não $\argmax_a \Qh(\raiz,a)$ — visitas são média sobre o histórico inteiro, robustas a um retorno afortunado.
    \item \textbf{Reaproveitamento da árvore.} Depois de agir, mantenha a subárvore do novo estado; descarte o resto.
    \item \textbf{Dois jogadores.} Troque o sinal do retorno a cada nível — senão joga bem \emph{contra si mesmo}.
    \item \textbf{Transposições.} Estados por ordens diferentes podem compartilhar estatísticas (grafo em vez de árvore).
    \item \textbf{Normalização de $\ret$.} Sem ela, $c$ perde significado — o bônus vira ruído ou tirania.
  \end{itemize}
\end{frame}

%% ------------------------------------------------------------
\begin{frame}{A árvore que o UCT constrói}
  \centering
  \begin{tikzpicture}[level distance=10mm, scale=0.95, transform shape,
      every node/.style={circle, draw=rlAzul, fill=rlAzul!10, thick,
                         minimum size=4mm, inner sep=0pt},
      level 1/.style={sibling distance=26mm},
      level 2/.style={sibling distance=11mm},
      level 3/.style={sibling distance=6mm},
      edge from parent/.style={draw, rlCinza, thick}]
    \node[fill=rlLaranja!25, draw=rlLaranja, minimum size=5mm] {}
      child {node {}
        child {node {}}
        child {node {}}}
      child {node[fill=rlLaranja!18, draw=rlLaranja] {}
        child {node[fill=rlLaranja!18, draw=rlLaranja] {}
          child {node {}}
          child {node[fill=rlLaranja!18, draw=rlLaranja] {}
            child {node {}} child {node {}}}
          child {node {}}}
        child {node {}
          child {node {}}}
        child {node {}}}
      child {node {}};
  \end{tikzpicture}

  \vfill
  \raggedright
  \begin{itemize}
    \item Laranja: o caminho principal (\emph{principal variation}) — onde a busca gastou quase todo o orçamento.
    \item Ruins visitados uma ou duas vezes e abandonados; promissores aprofundados. \textbf{Busca best-first altamente seletiva.}
  \end{itemize}
\end{frame}

%% ------------------------------------------------------------
\begin{frame}{O que se pode provar}
  \begin{teobox}{Consistência do UCT (Kocsis e Szepesvári, 2006)}
    Sob condições adequadas (retornos limitados, exploração suficiente), a
    probabilidade de o UCT escolher uma ação subótima na raiz tende a zero
    quando o número de simulações $K \to \infty$; a estimativa na raiz converge
    para o valor ótimo do sub-MDP.
  \end{teobox}

  \pause
  \begin{alertabox}
    Leia a letra miúda. Garantia \emph{assintótica}; pior caso \textbf{duplamente exponencial na profundidade} (Coquelin e Munos, 2007) — ``agulha no palheiro'' é catastroficamente lento.
  \end{alertabox}

  \pause
  \begin{itemize}
    \item Teoria: \emph{eventualmente} funciona. Prática: domínio benigno \emph{e} priors bons (Seção 3) encurtam drasticamente a busca.
  \end{itemize}
\end{frame}

%% ------------------------------------------------------------
\begin{frame}{Vantagens do MCTS}
  \begin{columns}[T]
    \begin{column}{0.5\textwidth}
      \begin{itemize}
        \item \textbf{Best-first seletivo} --- gasta onde importa.
        \item \textbf{Avaliação dinâmica} --- valor do estado \emph{atual}, não de todos.
        \item \textbf{Amostragem} quebra a maldição da dimensionalidade.
        \item \textbf{Modelo caixa-preta} --- só precisa de amostras.
      \end{itemize}
    \end{column}
    \begin{column}{0.5\textwidth}
      \begin{itemize}
        \item \textbf{\emph{Anytime}} --- interrompa quando quiser.
        \item \textbf{Paralelizável} --- simulações quase independentes (\emph{virtual loss}).
        \item \textbf{Sem heurística de domínio} --- dispensa a avaliação à mão do alfa-beta.
        \item \textbf{Melhora com computação} --- mais tempo, melhor jogada, sem retreinar nada.
      \end{itemize}
    \end{column}
  \end{columns}

  \vspace{0.2em}
  \begin{ideiabox}
    A última é a mais subestimada: computação em tempo de decisão → \emph{qualidade de decisão} — a propriedade buscada em modelos de raciocínio.
  \end{ideiabox}
\end{frame}

%% ------------------------------------------------------------
\begin{frame}{Limitações --- e o que se faz a respeito}
  \centering
  {\small
  \begin{tabular}{@{}p{4.3cm}p{7.4cm}@{}}
    \toprule
    \textbf{Limitação} & \textbf{Remédio usual} \\
    \midrule
    Ramificação enorme &
      \emph{prior} de política que concentra a busca (PUCT, Seção 3) \\
    Ações contínuas &
      \emph{progressive widening}: filhos crescem como $N^{\alpha}$ \\
    \emph{Rollouts} aleatórios são fracos &
      substituí-los por uma rede de valor $v_\theta$ \\
    Poucos dados por nó &
      compartilhar estatísticas entre nós (RAVE / AMAF) \\[-0.1em]
    Modelo aprendido, imperfeito &
      horizonte curto + valor aprendido, ou aprender o modelo no espaço de
      \emph{valores} (MuZero) \\
    \bottomrule
  \end{tabular}}

  \vspace{-4pt}
  \raggedright
  \begin{alertabox}
    MCTS pressupõe um simulador \emph{fiel} e \emph{barato}. Onde uma condição falha --- robótica, saúde, educação --- ele deixa de ser a ferramenta óbvia.
  \end{alertabox}
\end{frame}

%% ------------------------------------------------------------
\begin{frame}{Para discutir}
  \begin{quizbox}{}
    O UCT devolve $\argmax_a \Nv{\raiz,a}$ em vez de $\argmax_a \Qh(\raiz,a)$.
    Construa um exemplo pequeno em que as duas regras discordam. Qual das duas
    você confiaria mais, e por quê?
  \end{quizbox}
  \paradiscussao

  \pause
  \vfill
  \begin{quizbox}{}
    MCTS é ``eficiente em computação'' mas cada decisão exige milhares de
    chamadas ao simulador. Em que sentido, então, ele é \emph{eficiente em
    amostras} --- e em que sentido não é? Como isso se compara com o que vimos
    nas Aulas 11--12?
  \end{quizbox}
  \paradiscussao
\end{frame}

\section{De AlphaGo a MuZero}

%% ------------------------------------------------------------
\begin{frame}{Por que o \emph{Go} era o problema-símbolo}
  \begin{columns}[T]
    \begin{column}{0.54\textwidth}
      \begin{itemize}
        \item Tabuleiro $19\times19$; cerca de $10^{170}$ posições legais --- mais que átomos no universo observável.
        \item Ramificação $b\approx 250$ contra $b\approx 35$ do xadrez.
        \item Sem boa função de avaliação à mão: padrões globais sutis (vida e morte, influência, \emph{ko}).
        \item Alfa-beta com heurística (Deep Blue, 1997) \textbf{não transferiu}.
      \end{itemize}
    \end{column}
    \begin{column}{0.44\textwidth}
      \begin{ideiabox}
        2006+: MCTS levou o \emph{Go} do amador fraco ao forte. Faltava \emph{avaliação de posição} — redes profundas forneceram.
      \end{ideiabox}
      \begin{alertabox}
        2014: especialistas previam nível profissional a ``uma década de distância''. Levou dois anos.
      \end{alertabox}
    \end{column}
  \end{columns}
\end{frame}

%% ------------------------------------------------------------
\begin{frame}{A linha do tempo}
  \centering
  \begin{tikzpicture}[node distance=0mm]
    \draw[rlAzul!25, line width=2.2pt] (0,0) -- (11.6,0);
    \foreach \x/\ano/\ev/\dy in {%
      0.4/2006/{UCT\\ \color{rlCinza}Kocsis e\\ Szepesvári}/1.15,
      2.6/2016/{AlphaGo\\ \color{rlCinza}vence Lee\\ Sedol 4--1}/-1.35,
      5.0/2017/{AlphaGo Zero\\ \color{rlCinza}sem dados\\ humanos}/1.15,
      7.4/2018/{AlphaZero\\ \color{rlCinza}xadrez, shogi\\ e Go}/-1.35,
      9.8/2020/{MuZero\\ \color{rlCinza}modelo\\ aprendido}/1.15}{%
      \fill[rlLaranja] (\x,0) circle (2.6pt);
      \draw[rlLaranja, thick] (\x,0) -- (\x,{\dy*0.42});
      \node[anchor=center, align=center, font=\scriptsize\bfseries, color=rlAzul,
            text width=2.1cm] at (\x,\dy) {\ano\\[1pt] \ev};}
  \end{tikzpicture}

  \vfill
  \begin{itemize}
    \item Fio condutor: \textbf{remoção sucessiva de conhecimento humano} — mestres, \emph{features}, \emph{rollouts}, regras.
    \item MuZero: um simulador de si mesmo, uma busca e um sinal de recompensa.
  \end{itemize}
\end{frame}

%% ------------------------------------------------------------
\begin{frame}{AlphaGo (2016): três redes e uma busca}
  \centering
  \begin{tikzpicture}[node distance=5mm]
    \node[caixa, draw=rlTeal, fill=rlTeal!8, text width=2.5cm, inner sep=3.5pt]
      (sl) {\textbf{Rede de política SL}\\ $p_\sigma(a\mid s)$\\[2pt]
            \color{rlCinza}30M posições de partidas humanas};
    \node[caixa, draw=rlAzulClaro, fill=rlAzulClaro!8, text width=2.5cm,
          right=8mm of sl, inner sep=3.5pt]
      (rl) {\textbf{Rede de política RL}\\ $p_\rho(a\mid s)$\\[2pt]
            \color{rlCinza}auto-jogo, gradiente de política};
    \node[caixa, draw=rlAzul, fill=rlAzul!8, text width=2.5cm,
          right=8mm of rl, inner sep=3.5pt]
      (v) {\textbf{Rede de valor}\\ $v_\theta(s)$\\[2pt]
           \color{rlCinza}prevê o vencedor a partir do auto-jogo};
    \node[caixa, draw=rlLaranja, fill=rlLaranja!10, text width=2.5cm,
          right=8mm of v, inner sep=3.5pt]
      (m) {\textbf{MCTS}\\[2pt] \color{rlCinza}$p_\sigma$ como \emph{prior},
           $v_\theta$ e \emph{rollouts} como avaliação};
    \draw[trans destac] (sl)--(rl); \draw[trans destac] (rl)--(v);
    \draw[trans destac] (v)--(m);
    \draw[trans, dashed] (sl.south) to[bend right=22] (m.south);
  \end{tikzpicture}

  \vfill
  \raggedright
  \begin{itemize}
    \item Avaliação de folha por \textbf{mistura}:
          $\;\Vh(\folha) = (1-\lambda)\,v_\theta(\folha) + \lambda\,z_{\text{rollout}}$ — as fontes erram diferente; a média erra menos.
    \item A rede \emph{SL} funcionava melhor como \emph{prior} da busca que a RL — é mais \emph{diversa}.
    \item Vitórias: Fan Hui 5--0 (2015) e Lee Sedol 4--1 (2016).
  \end{itemize}
\end{frame}

%% ------------------------------------------------------------
\begin{frame}{AlphaGo Zero (2017): o que foi \emph{removido}}
  \centering
  {\small
  \begin{tabular}{@{}lcc@{}}
    \toprule
    & \textbf{AlphaGo} & \textbf{AlphaGo Zero} \\
    \midrule
    Partidas humanas          & sim, 30M posições & \color{rlVerde!70!black}\bfseries não \\
    \emph{Features} de domínio & sim, dezenas     & \color{rlVerde!70!black}\bfseries só o tabuleiro \\
    \emph{Rollouts} na busca   & sim              & \color{rlVerde!70!black}\bfseries não \\
    Redes separadas            & três             & \color{rlVerde!70!black}\bfseries uma, com duas cabeças \\
    Arquitetura                & convolucional    & residual \\
    \bottomrule
  \end{tabular}}

  \vspace{0.2em}
  \raggedright
  \begin{defbox}{Rede de duas cabeças}
    $(\bm{p}, v) = f_\theta(s)$, com $\bm{p} \in \Delta(\gA)$ o \emph{prior} de
    política e $v \in [-1,1]$ o valor da posição.
  \end{defbox}

  \vspace{-0.3em}
  \begin{itemize}\setlength{\itemsep}{0.2ex}
    \item Tronco compartilhado: as tarefas \emph{regularizam} uma à outra.
    \item 100--0 contra a versão que venceu Lee Sedol, após três dias de treino a partir de pesos aleatórios.
  \end{itemize}
\end{frame}

%% ------------------------------------------------------------
\begin{frame}{PUCT: o \emph{prior} entra na regra de seleção}
  \vspace{-7pt}
  \begin{teobox}{Regra de seleção do AlphaGo Zero / AlphaZero}
    \[
      a = \argmax_{a}\;
      \mdestaque{rlAzulClaro}{Q(s,a)}
      +
      \mdestaque{rlLaranja}{c_{\text{puct}}\, P(s,a)\,
      \frac{\sqrt{\sum_b \Nv{s,b}}}{1+\Nv{s,a}}},
      \quad P(s,a)=p_a,\;\; Q=\tfrac{W}{N} .
    \]
  \end{teobox}

  \pause
  Três diferenças em relação ao UCT:
  \begin{itemize}
    \item o bônus é \textbf{multiplicado pelo \emph{prior}} $P(s,a)$: ações implausíveis quase não são exploradas;
    \item o decaimento $1/(1+\Nv{s,a})$ é mais agressivo que $1/\sqrt{\Nv{s,a}}$;
    \item não há $\ln$: cresce com $\sqrt{\cdot}$ no total de visitas.
  \end{itemize}

  \pause
  \begin{ideiabox}
    Com $b\approx 250$, é o \emph{prior} que torna a busca viável: ela olha a sério para meia dúzia de jogadas. A rede não substitui a busca --- poda-a.
  \end{ideiabox}
\end{frame}

%% ------------------------------------------------------------
\begin{frame}{A ideia central: a busca é um operador de melhoria de política}
  Ao fim de $K$ simulações na raiz, defina a \textbf{política da busca}
  \[
    \pol_{\text{MCTS}}(a \mid s) \;=\;
    \frac{\Nv{s,a}^{1/\tau}}{\sum_b \Nv{s,b}^{1/\tau}},
    \qquad \tau > 0 \;\text{(temperatura)}.
  \]

  \pause
  \begin{teobox}{O laço de AlphaZero}
    Empiricamente, $\pol_{\text{MCTS}}$ é \emph{mais forte} que o \emph{prior}
    $\bm{p}$ da rede que a gerou. Então: use $\pol_{\text{MCTS}}$ como
    \textbf{alvo de treinamento} para $\bm{p}$, e o resultado da partida $z$
    como alvo para $v$. A rede melhorada torna a próxima busca melhor, que
    produz alvos ainda melhores.
  \end{teobox}

  \pause
  \begin{itemize}
    \item \textbf{Iteração de política} (Aula 3): busca = passo de \emph{melhoria}; auto-jogo = \emph{avaliação}.
    \item \emph{Expert iteration} (Anthony, Tian e Barber, 2017): busca = ``especialista lento''; rede = ``aprendiz rápido'' que \emph{amortiza}.
  \end{itemize}
\end{frame}

%% ------------------------------------------------------------
\begin{frame}{O laço de auto-jogo}
  \centering
  \begin{tikzpicture}[node distance=0mm]
    \def\rr{2.35}
    \node[caixa, draw=rlLaranja, fill=rlLaranja!10, text width=2.6cm,
          inner sep=4pt] (busca) at (0,\rr)
      {\textbf{MCTS}\\ \color{rlCinza}guiada por $f_\theta$,
       produz $\pol_{\text{MCTS}}$};
    \node[caixa, draw=rlAzulClaro, fill=rlAzulClaro!8, text width=2.6cm,
          inner sep=4pt] (jogo) at (\rr+1.6,0)
      {\textbf{Auto-jogo}\\ \color{rlCinza}$a_t \sim \pol_{\text{MCTS}}$
       até o fim, resultado $z$};
    \node[caixa, draw=rlTeal, fill=rlTeal!8, text width=2.6cm,
          inner sep=4pt] (dados) at (0,-\rr)
      {\textbf{Buffer}\\ \color{rlCinza}triplas
       $\bigl(s_t,\ \pol_{\text{MCTS}}(\cdot\mid s_t),\ z\bigr)$};
    \node[caixa, draw=rlAzul, fill=rlAzul!8, text width=2.6cm,
          inner sep=4pt] (treino) at (-\rr-1.6,0)
      {\textbf{Treino de $f_\theta$}\\ \color{rlCinza}$\bm{p}\to\pol_{\text{MCTS}}$,
       $v \to z$};
    \draw[trans destac] (busca) to[bend left=22] (jogo);
    \draw[trans destac] (jogo) to[bend left=22] (dados);
    \draw[trans destac] (dados) to[bend left=22] (treino);
    \draw[trans destac] (treino) to[bend left=22] (busca);
    \node[font=\scriptsize\itshape, color=rlCinza, align=center] at (0,0)
      {nenhum dado\\ humano entra\\ neste ciclo};
  \end{tikzpicture}
\end{frame}

%% ------------------------------------------------------------
\begin{frame}{A perda de treinamento}
  \begin{teobox}{Objetivo do AlphaGo Zero / AlphaZero}
    Para cada posição $s$ amostrada do auto-jogo, com alvos
    $\bm{\pi} = \pol_{\text{MCTS}}(\cdot\mid s)$ e $z \in \{-1,0,+1\}$:
    \[
      \ell(\theta) \;=\;
      \termo{(z - v_\theta(s))^2}{regressão de valor}
      \;-\;
      \termo{\bm{\pi}^{\!\top} \log \bm{p}_\theta(s)}{destilação da busca}
      \;+\;
      \termo{c\,\norm{\theta}^2}{regularização} .
    \]
  \end{teobox}

  \pause
  \begin{itemize}
    \item 2º termo: \emph{entropia cruzada} — supervisionado, rótulo fabricado pela própria busca.
    \item 1º termo fecha o ciclo: sem $v$ bom, a busca degrada e os rótulos pioram.
    \item Sem gradiente de política, razão de importância ou região de confiança --- o problema difícil de RL foi empurrado para dentro da busca.
  \end{itemize}
\end{frame}

%% ------------------------------------------------------------
\begin{frame}{Duas engenhocas indispensáveis}
  \vspace{-9pt}
  \begin{defbox}{Ruído de Dirichlet na raiz}
    \[
      P(\raiz, a) \;=\; (1-\varepsilon)\, p_a \;+\; \varepsilon\, \eta_a,
      \qquad \bm{\eta} \sim \mathrm{Dir}(\alpha), \quad \varepsilon \approx 0{,}25 .
    \]
    Garante que toda jogada legal tenha alguma chance de ser explorada, mesmo
    que a rede a considere impossível. Sem isso, o auto-jogo colapsa em um
    repertório estreito.
  \end{defbox}

  \pause
  \begin{defbox}{Temperatura $\tau$}
    Nas primeiras jogadas, $\tau = 1$ (amostra proporcional às visitas), o que
    gera \textbf{diversidade de aberturas}; depois $\tau \to 0$
    (praticamente $\argmax$), jogo forte.
  \end{defbox}

  \pause
  \begin{alertabox}
    As duas são a versão de auto-jogo do problema de \emph{exploração} das Aulas 10--12: sem elas o sistema aprende a bater apenas em si mesmo.
  \end{alertabox}
\end{frame}

%% ------------------------------------------------------------
\begin{frame}{AlphaZero (2018): a mesma receita, três jogos}
  \begin{itemize}
    \item Um único algoritmo, sem conhecimento de domínio além das regras, do zero por auto-jogo em \textbf{Go}, \textbf{xadrez} e \textbf{shogi}.
    \item Derrotou o Stockfish (xadrez), o elmo (shogi) e a versão anterior do AlphaGo Zero.
    \item Dezenas de milhares de posições por jogada vs.\ dezenas de milhões do Stockfish --- e ainda assim jogava melhor.
  \end{itemize}

  \pause
  \begin{ideiabox}
    A lição não é ``busca'' nem ``rede'', é a \textbf{divisão de trabalho}: rede = julgamento posicional (largura); busca = cálculo concreto (profundidade).
  \end{ideiabox}

  \pause
  \begin{alertabox}
    Força vs.\ motores: depende de tempo por jogada, \emph{hardware} e livro de aberturas — conclusão qualitativa sólida, números nem tanto.
  \end{alertabox}
\end{frame}

%% ------------------------------------------------------------
\begin{frame}{MuZero (2020): e se nem as regras forem dadas?}
  \centering
  \begin{tikzpicture}[node distance=4mm]
    \node[caixa, draw=rlTeal, fill=rlTeal!8, text width=1.9cm, inner sep=3pt]
      (h) {$h_\theta$\\ \color{rlCinza}\tiny observação $\to$ estado latente $z_0$};
    \node[caixa, draw=rlAzulClaro, fill=rlAzulClaro!8, text width=1.9cm,
          right=9mm of h, inner sep=3pt]
      (g1) {$g_\theta$\\ \color{rlCinza}\tiny $(z_0,a_1)\to(z_1,\hat r_1)$};
    \node[caixa, draw=rlAzulClaro, fill=rlAzulClaro!8, text width=1.9cm,
          right=9mm of g1, inner sep=3pt]
      (g2) {$g_\theta$\\ \color{rlCinza}\tiny $(z_1,a_2)\to(z_2,\hat r_2)$};
    \node[caixa, draw=rlLaranja, fill=rlLaranja!10, text width=1.9cm,
          right=9mm of g2, inner sep=3pt]
      (f) {$f_\theta$\\ \color{rlCinza}\tiny $z_k \to (\bm{p}_k, v_k)$};
    \draw[trans destac] (h)--(g1); \draw[trans destac] (g1)--(g2);
    \draw[trans destac] (g2)--(f);
    \draw[trans] (h.south) to[bend right=26]
      node[below, font=\tiny, color=rlCinza, pos=0.55]
      {$f_\theta$ é aplicada a \emph{todo} $z_k$} (f.south);
  \end{tikzpicture}

  \vfill
  \raggedright
  \begin{itemize}
    \item Sem decodificador: $z_k$ \textbf{não} reconstrói o estado real — basta prever \emph{recompensa}, \emph{valor} e \emph{política}, o que a busca consome.
    \item Busca inteiramente no espaço latente; modelo aprendido \emph{para servir ao planejamento}, não para ser fiel.
    \item Igualou AlphaZero nos jogos de tabuleiro; estado da arte em Atari, onde ninguém tem as regras.
  \end{itemize}
\end{frame}

%% ------------------------------------------------------------
\begin{frame}{Para discutir}
  \begin{quizbox}{}
    O alvo de política é a distribuição de visitas $\bm{\pi}$, e o alvo de valor
    é o resultado final $z$ da partida. Por que não usar $\Qh(\raiz, a)$ da busca
    como alvo de valor, já que ele parece mais informativo?
  \end{quizbox}
  \paradiscussao

  \pause
  \vfill
  \vspace{-2.5pt}
  \begin{quizbox}{}
    O auto-jogo em jogos de dois jogadores de soma zero gera automaticamente um
    adversário do próprio nível --- um currículo. Que ingrediente equivalente
    existiria em um problema de um jogador só, como controle de um robô?
  \end{quizbox}
  \paradiscussao
\end{frame}

%% ------------------------------------------------------------
\begin{frame}{O que generaliza --- e o que não}
  \begin{columns}[T]
    \begin{column}{0.49\textwidth}
      \textbf{\color{rlVerde!70!black}Transfere bem}
      \begin{itemize}\setlength{\itemsep}{0.2ex}
        \item Busca como operador de melhoria; rede como amortização da busca.
        \item \emph{Prior} de política para podar ramificação alta.
        \item Trocar computação de decisão por qualidade de decisão.
        \item Aprender o modelo no espaço do que o planejador consome.
      \end{itemize}
    \end{column}
    \begin{column}{0.49\textwidth}
      \textbf{\color{rlVermelho}Não transfere de graça}
      \begin{itemize}\setlength{\itemsep}{0.2ex}
        \item Simulador perfeito, determinístico e barato.
        \item Recompensa terminal não ambígua ($\pm1$).
        \item Auto-jogo como fonte gratuita de currículo.
        \item Estados episódicos e reinicializáveis à vontade.
      \end{itemize}
    \end{column}
  \end{columns}

  \vspace{-2pt}
  \begin{alertabox}
    Fora de jogos o gargalo raramente é a busca: é um modelo confiável por vinte passos e uma recompensa que signifique o que queremos.
  \end{alertabox}
\end{frame}

%% ------------------------------------------------------------
\begin{frame}{MCTS além de jogos}
  \begin{itemize}
    \item \textbf{Planejamento e controle} --- MPC amostral, POMDPs (POMCP, DESPOT), escalonamento, roteamento, síntese de programas.
    \item \textbf{Química e matemática} --- retrossíntese, descoberta de algoritmos (AlphaTensor), demonstração de teoremas (AlphaProof).
    \item \textbf{Modelos de linguagem} --- busca sobre passos de raciocínio; modelo como \emph{prior}, verificador no papel da rede de valor. ``Gerar, avaliar, destilar de volta'' é o laço de AlphaZero.
  \end{itemize}

  \pause
  \begin{alertabox}
    A analogia tem limite: ``estado'' não-markoviano, ramo essencialmente infinito, sem regra de vitória. A peça mais difícil é o \emph{verificador}, não a busca.
  \end{alertabox}
\end{frame}

\section{Fechamento}

%% ------------------------------------------------------------
\begin{frame}{Resumo da aula}
  \begin{enumerate}
    \item \textbf{Planejamento em tempo de decisão.} Com um simulador, vale gastar computação local para decidir melhor \emph{agora}.
    \item \textbf{Busca de Monte Carlo simples} $=$ um passo de melhoria de política; expectimax custaria $(\abs{\gS}\abs{\gA})^H$.
    \item \textbf{MCTS} escapa amostrando (largura) e truncando (profundidade); árvore assimétrica, \emph{anytime}.
    \item \textbf{UCT} trata cada nó como um bandit; converge, com garantias fracas no pior caso.
    \item \textbf{AlphaZero}: \emph{rollouts} $\to$ $v_\theta$; busca guiada por $\bm{p}_\theta$ via PUCT; treino nos alvos que a própria busca produz — iteração de política.
    \item \textbf{MuZero} aprende também o modelo, no espaço do que o planejador consome.
  \end{enumerate}
\end{frame}

%% ------------------------------------------------------------
\begin{frame}{Formulário}
  \begin{columns}[T]
    \begin{column}{0.5\textwidth}
      {\small
      \textbf{UCT}
      \[ a = \argmax_a \Qh(i,a) + c\sqrt{\tfrac{\ln \Nv{i}}{\Nv{i,a}}} \]
      \textbf{PUCT}
      \[ a = \argmax_a Q(s,a) + c_{\text{puct}} P(s,a)\tfrac{\sqrt{\sum_b \Nv{s,b}}}{1+\Nv{s,a}} \]
      \textbf{Política da busca}
      \[ \pol_{\text{MCTS}}(a\mid s) \propto \Nv{s,a}^{1/\tau} \]}
    \end{column}
    \begin{column}{0.5\textwidth}
      {\small
      \textbf{Retropropagação}
      \[ \ret \gets r + \gamma \ret, \quad
         \Qh \gets \tfrac{W + \ret}{N+1} \]
      \textbf{Perda AlphaZero}
      \[ \ell = (z-v)^2 - \bm{\pi}^{\!\top}\log \bm{p} + c\norm{\theta}^2 \]
      \textbf{Ruído na raiz}
      \[ P \gets (1-\varepsilon)\bm{p} + \varepsilon\,\bm{\eta},\;
         \bm{\eta}\sim\mathrm{Dir}(\alpha) \]}
    \end{column}
  \end{columns}
\end{frame}

%% ------------------------------------------------------------
\begin{frame}{Erros comuns}
  \begin{alertabox}
    \begin{itemize}\setlength{\itemsep}{0.3ex}
      \item Devolver $\argmax_a \Qh$ em vez de $\argmax_a \Nv{\cdot,a}$ --- traído por um único \emph{rollout} sortudo.
      \item Esquecer de inverter o sinal do retorno a cada nível em jogos de dois jogadores.
      \item Usar $c$ do UCT sem normalizar os retornos.
      \item Descartar a árvore inteira a cada jogada.
      \item Treinar $\bm{p}$ contra o $\argmax$ da busca em vez da distribuição de visitas --- perde-se a informação de incerteza.
      \item Supor que um modelo com $1\%$ de erro por passo é utilizável em um horizonte de $50$ passos.
    \end{itemize}
  \end{alertabox}
\end{frame}

%% ------------------------------------------------------------
\begin{frame}{Leituras}
  {\small
  \begin{itemize}\setlength{\itemsep}{0.25ex}
    \item L. Kocsis e C. Szepesvári. Bandit based Monte-Carlo planning.
          \emph{ECML}, 2006. \hfill \color{rlCinza}\scriptsize o artigo do UCT
    \item \color{black} C. Browne et al. A survey of Monte Carlo tree search
          methods. \emph{IEEE TCIAIG}, 4(1):1--43, 2012.
    \item P.-A. Coquelin e R. Munos. Bandit algorithms for tree search.
          \emph{UAI}, 2007. \hfill \color{rlCinza}\scriptsize os limites do UCT
    \item \color{black} D. Silver et al. Mastering the game of Go with deep
          neural networks and tree search. \emph{Nature}, 529:484--489, 2016.
    \item D. Silver et al. Mastering the game of Go without human knowledge.
          \emph{Nature}, 550:354--359, 2017.
    \item D. Silver et al. A general reinforcement learning algorithm that
          masters chess, shogi, and Go through self-play. \emph{Science},
          362:1140--1144, 2018.
    \item T. Anthony, Z. Tian e D. Barber. Thinking fast and slow with deep
          learning and tree search. \emph{NeurIPS}, 2017.
    \item J. Schrittwieser et al. Mastering Atari, Go, chess and shogi by
          planning with a learned model. \emph{Nature}, 588:604--609, 2020.
    \item R. Sutton e A. Barto. \emph{Reinforcement Learning}, 2ª ed., Cap.\ 8.
    \item E. Brunskill. \emph{CS234: Reinforcement Learning}, Stanford ---
          \emph{Lecture 13: Monte Carlo Tree Search}.
  \end{itemize}}
\end{frame}

%% ------------------------------------------------------------
\begin{frame}{Próxima aula}
  \begin{itemize}
    \item \textbf{Recompensas em RL}: de onde vem $\rec$; o que acontece quando está errada; como especificá-la sem comportamento indesejado.
    \item A pergunta que fica: MCTS otimiza \emph{muito bem} o que lhe for dado como recompensa — virtude apenas se a recompensa estiver certa.
  \end{itemize}

  \vfill
  \centering
  {\color{rlCinza}\small Material adaptado e expandido de \textbf{CS234:
  Reinforcement Learning}, Emma Brunskill, Stanford University.\\
  Prof.\ Marcelo Keese Albertini --- Faculdade de Computação, UFU.}
\end{frame}

\end{document}
